\section*{الفصل \ref{ch:g9:maps-of-the-earth} --- خرائط الأرض كلها}

\begin{solution}{exo:g9:maps-of-the-earth:1}
(أ) \omterm{def:g9:maps-of-the-earth:projection}{الإسقاط الخرائطي}. (ب) \omterm{def:g9:maps-of-the-earth:contour}{خط الكنتور}. (ج) \omterm{def:g9:maps-of-the-earth:numerical-scale}{المقياس العددي}. (د) \omterm{def:g9:maps-of-the-earth:contour}{الفاصل
الكنتوري}. (ه) \omterm{def:g9:maps-of-the-earth:remote-sensing}{الاستشعار عن بعد}.
\end{solution}

\begin{solution}{exo:g9:maps-of-the-earth:2}
(أ) $4 \times 25\,000 = \num{100000}$~cm $= \qty{1}{km}$. (ب) $3 \times
50\,000 = \num{150000}$~cm $= \qty{1.5}{km}$. (ج) $2.5 \times 1\,000\,000
= \num{2500000}$~cm $= \qty{25}{km}$. (د) $\qty{1}{km} = \num{100000}$~cm؛
$100\,000 \div 25\,000 = \qty{4}{cm}$.
\end{solution}

\begin{solution}{exo:g9:maps-of-the-earth:3}
(أ) خطأ: لا يمكن بسط الكرة دون تشويه الأطوال، فكل \omterm{def:g3:continents-and-oceans:world-map}{خريطة للعالم} تشوّه
شيئًا ما. (ب) صواب. (ج) خطأ: \omterm{def:g2:drawing-to-map:map}{الخريطة} المتساوية المساحة تحفظ المساحات،
لكنها تثني الأشكال، ولا سيما قرب أطرافها. (د) صواب.
\end{solution}

\begin{solution}{exo:g9:maps-of-the-earth:4}
(أ) الانحدار شديد. (ب) الأرض لطيفة الانحدار. (ج) \omterm{def:g1:landscapes-around-us:landform}{تلة}، قمتها داخل أصغر
خط. (د) لنقطة من الأرض \omterm{def:g4:world-physical-map:altitude}{ارتفاع} واحد فقط، فلا يمكن أن تقع على خطين
بارتفاعين مختلفين.
\end{solution}

\begin{solution}{exo:g9:maps-of-the-earth:5}
(أ) عند \ang{60} شمالًا تُرسم \omterm{def:g1:landscapes-around-us:water}{البحيرة} أكبر بأربع مرات مما ينبغي: فتبدو بحجم
$4 \times 1000 = \qty{4000}{km^2}$. وعند \omterm{def:g3:continents-and-oceans:equator}{خط الاستواء} لا تُكبَّر المساحات، لذلك
\omterm{def:g1:landscapes-around-us:water}{فالبحيرة} المرسومة بالحجم نفسه مساحتها الحقيقية \qty{4000}{km^2}.
(ب) $4 \div 1.33 \approx 3$: أي أكبر بثلاث مرات. (ج) المساحات، والمسافات
من الشمال إلى الجنوب، تكبر بلا حد نحو \omterm{def:g3:continents-and-oceans:poles}{القطبين}: فيكون القطبان على بعد لا
نهائي من \omterm{def:g3:continents-and-oceans:equator}{خط الاستواء}.
\end{solution}

\begin{solution}{exo:g9:maps-of-the-earth:6}
(أ) \qty{50}{m}. (ب) القمة على \omterm{def:g4:world-physical-map:altitude}{ارتفاع} \qty{544}{m}؛ \omterm{def:g9:maps-of-the-earth:contour}{وخط الكنتور} الرئيسي
\omterm{def:g2:water-on-earth:ocean}{المحيط} بها هو خط \qty{500}{m}. (ج) في الغرب: \omterm{def:g1:landscapes-around-us:water}{فالنهر} يجري منحدرًا، لذلك
يكون قاعه أخفض في اتجاه المصب، أي في الشرق. (د) بين خطي 300
و\qty{400}{m} تقريبًا: فهناك يكون الانحدار أشد، وهو ألطف قرب قاع \omterm{def:g1:landscapes-around-us:landform}{الوادي}
وقرب القمة.
\end{solution}

\begin{solution}{exo:g9:maps-of-the-earth:7}
(أ) C على \omterm{def:g4:world-physical-map:altitude}{ارتفاع} نحو \qty{400}{m}، وD على \omterm{def:g4:world-physical-map:altitude}{ارتفاع} نحو \qty{530}{m}.
(ب) نحو \qty{135}{m}، على بعد نحو \qty{1.3}{km} من C، في قاع \omterm{def:g1:landscapes-around-us:landform}{الوادي}.
(ج) نحو 0.3 إلى \qty{0.4}{km}. (د) الارتفاعات صغيرة مقارنةً بالمسافات:
فدون المط، تبدو \omterm{def:g1:landscapes-around-us:landform}{التلال} شبه مستوية.
\end{solution}

\begin{solution}{exo:g9:maps-of-the-earth:8}
(أ) $2 \times 25\,000 = \num{50000}$~cm $= \qty{500}{m}$. (ب) $300 - 200 =
\qty{100}{m}$. (ج) $100 \div 500 = 0.2$، أي انحدار نسبته \qty{20}{\%}. (د) ألطف:
إذ تساوي \qty{4}{cm} مسافة \qty{1000}{m}، فالانحدار $100 \div 1000 = \qty{10}{\%}$،
أي نصف الانحدار السابق.
\end{solution}

\begin{solution}{exo:g9:maps-of-the-earth:9}
(أ) $14 \times 18 = 252$ مدارًا. (ب) $14 \times 103 = 1442$ دقيقة، أي نحو
\qty{24}{h}: يوم كامل. (ج) $5.5 \div 1 = 5.5$ مرة. (د) مثلًا: \omterm{def:g8:ice-on-earth:glacier}{نهر جليدي}
يتراجع، أو غابة تُقطع، أو بلدة تكبر، أو \omterm{def:g1:landscapes-around-us:water}{ساحل} أو \omterm{def:g1:landscapes-around-us:water}{نهر}
يتحرك.
\end{solution}

\begin{solution}{exo:g9:maps-of-the-earth:10}
(أ) $20\,200 \div 300\,000 \approx 0.067$~s $= 6.7 \times 10^{-2}$~s.
(ب) $300\,000 \times 10^{-3} = \qty{300}{km}$. (ج) خطأ ضئيل في الزمن
يُحدث خطأ هائلًا في المسافة: فيجب أن تكون الساعات دقيقة للغاية.
\end{solution}

\begin{solution}{exo:g9:maps-of-the-earth:11}
(أ) المسافة إلى A: $\sqrt{3.6^2 + 4.8^2} = \sqrt{12.96 + 23.04} =
\sqrt{36} = 6$؛ والمسافة إلى B: $\sqrt{6.4^2 + 4.8^2} = \sqrt{40.96 +
23.04} = \sqrt{64} = 8$. (ب) $(3.6, -4.8)$، وهي صورتها المرآتية بالنسبة
إلى المستقيم AB. (ج) المسافة من C إلى $(3.6, 4.8)$: $\sqrt{3.6^2 + 5.2^2} =
\sqrt{40} \approx 6.3$؛ وإلى $(3.6, -4.8)$: $\sqrt{3.6^2 + 14.8^2} =
\sqrt{232} \approx 15.2$. فالمستقبِل عند $(3.6, 4.8)$. (د) مسافة واحدة
تعطي دائرة (كرة في الفضاء)، ومسافتان تعطيان نقطتين ممكنتين؛ ويلزم مسافة
أخرى للاختيار.
\end{solution}

\begin{solution}{exo:g9:maps-of-the-earth:12}
(أ) $29 \div 2.1 \approx 14$ مرة. (ب) $2.1 \times 10.5 \approx
22$~مليون~\unit{km^2}. (ج) تبدو غرينلاند كأنها تغطي نحو 22~مليون
\unit{km^2}، وهذا قريب من 29~مليونًا لإفريقيا: ففي \omterm{def:g2:drawing-to-map:map}{خريطة} الجدار \omterm{def:g9:maps-of-the-earth:projection}{بإسقاط
مركاتور}، تبدو الاثنتان بالحجم نفسه تقريبًا، مع أن مساحة إفريقيا 14 ضعف
مساحتها. (د) \omterm{def:g9:maps-of-the-earth:projection}{إسقاط متساوي المساحة}.
\end{solution}

\begin{solution}{pb:g9:maps-of-the-earth:1}
\textbf{1.} $50\,000$~cm $= \qty{500}{m}$.

\textbf{2.} $3.6 \times 500 = \qty{1800}{m} = \qty{1.8}{km}$.

\textbf{3.} B--W: $2.5 \times 500 = \qty{1250}{m}$؛ W--T: $2.35 \times 500
\approx \qty{1180}{m}$.

\textbf{4.} نحو الشمال الشرقي.

\textbf{5.} \qty{50}{m}.

\textbf{6.} المرحلة S--B تنزل: $296 - 143 = \qty{153}{m}$.

\textbf{7.} B--W: $335 - 143 = \qty{192}{m}$؛ W--T: $544 - 335 =
\qty{209}{m}$.

\textbf{8.} $209 \div 1180 \approx 0.18$، أي نحو \qty{18}{\%}.

\textbf{9.} الخط المستقيم من S إلى T يعبر \omterm{def:g1:landscapes-around-us:water}{النهر} حيث لا جسر؛ فالدرب ينزل
إلى الجسر عند B.

\textbf{10.} $1.8 + 1.25 + 1.18 \approx \qty{4.2}{km}$ بسرعة \qty{4}{km} في
الساعة: نحو 63 دقيقة؛ والصعود $192 + 209 = \qty{401}{m}$: نحو 40 دقيقة
إضافية. والمجموع نحو ساعة و43 دقيقة.

\textbf{11.} من T تنزل إلى B، ثم تصعد من B إلى S: $296 -
143 = \qty{153}{m}$.

\textbf{12.} الطول: $1.8 + 1.25 + 1.18 \approx \textbf{4.2 km}$؛ ومجموع
الصعود: $192 + 209 = \textbf{401 m}$.
\end{solution}
